% =====================================================================
%  1. OBJETIVO
% =====================================================================
\section{Objetivo}
Implementar en \matlab{} una funci\'{o}n con la signatura
\begin{center}
  \texttt{[coste, ruta] = dijkstra(G, origen, destino)}
\end{center}
que devuelva el coste y la secuencia de nodos del camino \'{o}ptimo entre dos nodos de un grafo dirigido y ponderado, y validarla sobre los dos grafos del enunciado: el grafo de ejemplo de $7$ nodos y la matriz \texttt{J} de \texttt{grafos.mat} ($25$ nodos, dirigida y asim\'{e}trica).

El convenio del enunciado es que $G(i,j) = 0$ representa ausencia de arco $i\to j$ y los pesos son no negativos, condici\'{o}n necesaria para que Dijkstra sea correcto.


% =====================================================================
%  2. IMPLEMENTACIÓN
% =====================================================================
\section{Implementaci\'{o}n}

La funci\'{o}n sigue el esquema cl\'{a}sico de Dijkstra (selecci\'{o}n del nodo abierto con menor distancia, relajaci\'{o}n de aristas, marcado como cerrado), por lo que aqu\'{i} solo comentamos las decisiones que no son evidentes del enunciado.

\paragraph{Selecci\'{o}n del m\'{i}nimo por b\'{u}squeda lineal.} Para los tama\~{n}os con los que se trabaja ($N \leq 25$), una b\'{u}squeda lineal sobre los nodos abiertos en cada iteraci\'{o}n es m\'{a}s clara y r\'{a}pida que mantener una cola de prioridad. La complejidad $O(N^2)$ resultante es perfectamente aceptable para este problema; con grafos mucho m\'{a}s grandes habr\'{i}a que reconsiderarlo.

\paragraph{Aceptar el \texttt{struct} de \texttt{load} directamente.} El enunciado entrega los grafos dentro de \texttt{grafos.mat}. Para evitar el paso intermedio de extraer la matriz a mano, la funci\'{o}n acepta tanto una matriz como el \texttt{struct} devuelto por \texttt{load}: si recibe el \texttt{struct}, busca el campo \texttt{J} (la matriz del enunciado) y, en su defecto, la primera matriz cuadrada num\'{e}rica que encuentre. As\'{i} se puede invocar como \texttt{dijkstra(load('grafos.mat'), 1, 7)} sin pasos intermedios.

\paragraph{Corte temprano.} Como interesa el camino a un destino concreto, el bucle se detiene en cuanto se extrae $u = t$ del conjunto abierto. No cambia la complejidad asint\'{o}tica, pero reduce el trabajo medio.

\paragraph{Validaciones de entrada.} Antes del bucle se comprueba que \texttt{G} es cuadrada, que los costes son finitos y no negativos, y que los \'{i}ndices de origen y destino est\'{a}n dentro del rango. Si \texttt{origen == destino} se devuelve \texttt{coste = 0} sin entrar al bucle. Si la diagonal contiene autolazos, se avisa pero no se aborta: el filtro $G(u,v) > 0$ al recorrer vecinos los ignora autom\'{a}ticamente.

\paragraph{Reconstrucci\'{o}n de la ruta.} Si $d[t] = \infty$ al terminar, se devuelve \texttt{coste = Inf} y \texttt{ruta = []} (destino inalcanzable). En caso contrario se reconstruye la ruta retrocediendo por el vector de predecesores hasta el origen.


% =====================================================================
%  3. PRUEBAS
% =====================================================================
\section{Pruebas}

\subsection{Grafo de ejemplo (7 nodos)}
Sobre el grafo $G$ de $7$ nodos del enunciado:
\begin{equation*}
  G = \begin{bmatrix}
    0 & 2 & 0 & 0 & 9 & 0 & 0 \\
    2 & 0 & 1 & 2 & 0 & 0 & 0 \\
    0 & 1 & 0 & 5 & 0 & 0 & 0 \\
    0 & 2 & 5 & 0 & 1 & 2 & 4 \\
    9 & 0 & 0 & 1 & 0 & 4 & 0 \\
    0 & 0 & 0 & 2 & 4 & 0 & 1 \\
    0 & 0 & 0 & 4 & 0 & 1 & 0
  \end{bmatrix}
\end{equation*}

\noindent la llamada \texttt{dijkstra(G, 1, 7)} devuelve \texttt{coste = 7} y \texttt{ruta = [1 2 4 6 7]}, que coincide con el resultado esperado del enunciado.

\begin{figure}[H]
  \centering
  \includegraphics[width=0.7\linewidth]{dijkstra_ejemplo7}
  \caption{Grafo de $7$ nodos con la ruta \'{o}ptima de $1$ a $7$ resaltada en rojo. Coste $7$, ruta $[1\,2\,4\,6\,7]$.}
  \label{fig:dijkstra_ejemplo7}
\end{figure}

\subsection{Matriz J de grafos.mat (25 nodos)}
La matriz \texttt{J} es un grafo dirigido y asim\'{e}trico de $25$ nodos. La asimetr\'{i}a es relevante: por ejemplo $J(1,7) \neq J(7,1)$, lo que hace que las rutas de ida y vuelta entre el mismo par de nodos no sean iguales.

\begin{table}[H]
  \centering
  \caption{Resultados de \texttt{dijkstra} sobre la matriz \texttt{J} ($25$ nodos, dirigida).}
  \label{tab:dijkstra_validacion}
  \begin{tabularx}{0.98\linewidth}{c c c X}
    \toprule
    \textbf{Origen} & \textbf{Destino} & \textbf{Coste} & \textbf{Ruta} \\
    \midrule
    1  & 19 & $16$  & $[1\,7\,13\,17\,23\,18\,19]$ \\
    19 & 1  & $19$  & $[19\,20\,15\,10\,5\,9\,8\,3\,2\,1]$ \\
    25 & 19 & $11$  & $[25\,20\,15\,14\,18\,19]$ \\
    19 & 25 & $4$   & $[19\,20\,25]$ \\
    1  & 23 & $12$  & $[1\,7\,13\,17\,23]$ \\
    23 & 1  & $22$  & $[23\,18\,14\,15\,10\,5\,9\,8\,3\,2\,1]$ \\
    1  & 24 & $\infty$ & no alcanzable \\
    23 & 22 & $40$  & $[23\,18\,14\,15\,10\,5\,9\,8\,3\,2\,1\,7\,13\,17\,12\,11\,16\,22]$ \\
    \bottomrule
  \end{tabularx}
\end{table}

Conviene fijarse en dos casos. Para $1 \to 23$ el coste es $12$ y la ruta usa $5$ nodos; para $23 \to 1$ el coste sube a $22$ y la ruta a $10$ nodos. M\'{a}s extremo es $23 \to 22$: aunque $22$ y $23$ son nodos pr\'{o}ximos en el mapa, los arcos disponibles obligan a recorrer casi todo el grafo (coste $40$, $18$ nodos). El caso $1 \to 24$ confirma que la funci\'{o}n maneja correctamente la inalcanzabilidad.

Tambi\'{e}n se han probado los casos l\'{i}mite \texttt{origen == destino} (devuelve coste $0$ y ruta de un solo nodo) y entrada con pesos negativos (rechazada con error en la validaci\'{o}n).


% =====================================================================
%  4. COMENTARIOS FINALES
% =====================================================================
\section{Comentarios finales}

La implementaci\'{o}n cumple lo que pide el enunciado y se ha validado contra todos los casos propuestos. La b\'{u}squeda lineal del m\'{i}nimo no es la opci\'{o}n m\'{a}s eficiente posible, pero para los tama\~{n}os de \texttt{grafos.mat} no merece la pena complicarla con una cola de prioridad; en grafos mucho m\'{a}s dispersos y grandes ser\'{i}a el siguiente paso a mirar.

Esta misma funci\'{o}n se reutiliza en Lab-MR 5 como referencia para comparar contra A*, y en Lab-MR 6 como planificador global cuya ruta se entrega luego a la navegaci\'{o}n local reactiva.
